\pdfminorversion=6
\usepackage[letterpaper,margin=1in]{geometry}
\usepackage[T1]{fontenc}
\usepackage[utf8]{inputenc}
\usepackage{lmodern,microtype}
\usepackage{amsmath,amssymb,mathtools,amsthm}
\usepackage{bm}
\usepackage[normalem]{ulem}
\usepackage{xcolor,graphicx,enumitem,dsfont,tikz}
\usetikzlibrary{arrows.meta,calc}

\usepackage[titles]{tocloft}
\renewcommand{\cftsecleader}{\cftdotfill{\cftdotsep}}
\renewcommand{\cftdotsep}{2}
\setlength{\cftbeforesecskip}{6pt}

\usepackage{hyperref}
\hypersetup{
  colorlinks=true,
  linkcolor=blue!60!black,
  citecolor=green!40!black,
  urlcolor=blue!70!black,
  linktoc=all,
  bookmarksdepth=2,
  pdftitle={The limit shape of uniform random antichains in [0,1] to the power d},
  pdfsubject={Limit shapes, permutation limits, and applications to coordinate orders}
}

\usepackage[nameinlink,capitalize,noabbrev]{cleveref}
\crefname{section}{Section}{Sections}
\Crefname{section}{Section}{Sections}
\crefname{subsection}{Section}{Sections}
\Crefname{subsection}{Section}{Sections}
\crefname{subsubsection}{Section}{Sections}
\Crefname{subsubsection}{Section}{Sections}
\renewcommand*{\sectionautorefname}{Section}
\renewcommand*{\subsectionautorefname}{Section}
\renewcommand*{\subsubsectionautorefname}{Section}

\usepackage[
  backend=biber,
  style=alphabetic,
  maxalphanames=3,
  maxnames=8,
  maxbibnames=20,
  giveninits=true
]{biblatex}
\renewcommand*{\multicitedelim}{\addcomma\space}
\addbibresource{references.bib}
\setlength{\bibitemsep}{2pt}
\AtBeginBibliography{\small\sloppy}

\setlength{\parskip}{3pt}
\setlength{\emergencystretch}{2em}
\setcounter{tocdepth}{2}
\setcounter{secnumdepth}{3}

\numberwithin{equation}{section}
\newtheorem{theorem}{Theorem}[section]

\newtheorem{proposition}[theorem]{Proposition}
\newtheorem{lemma}[theorem]{Lemma}
\newtheorem{corollary}[theorem]{Corollary}
\providecommand*{\propositionautorefname}{Proposition}
\providecommand*{\lemmaautorefname}{Lemma}
\providecommand*{\corollaryautorefname}{Corollary}

\theoremstyle{definition}
\newtheorem{question}[theorem]{Question}
\providecommand*{\questionautorefname}{Question}

\newtheorem{definition}[theorem]{Definition}
\newtheorem{conjecture}[theorem]{Conjecture}
\newtheorem{problem}[theorem]{Problem}
\newtheorem{assumption}[theorem]{Assumption}

\theoremstyle{remark}
\newtheorem{remark}[theorem]{Remark}

\crefname{theorem}{Theorem}{Theorems}
\crefname{proposition}{Proposition}{Propositions}
\crefname{lemma}{Lemma}{Lemmas}
\crefname{corollary}{Corollary}{Corollaries}
\crefname{conjecture}{Conjecture}{Conjectures}
\crefname{problem}{Problem}{Problems}
\crefname{question}{Question}{Questions}
\Crefname{question}{Question}{Questions}

\setlist{topsep=3pt,itemsep=2pt,parsep=0pt}



